\section{Voltage Follower}\label{sec:follower}
Mạch đệm (\textit{voltage follower}) có độ lợi bằng 1: $V_{OUT} = V_{(+)}$, trở kháng vào rất lớn, trở kháng ra rất nhỏ. Mạch dùng làm tầng đệm giữa nguồn có trở kháng ra lớn và tải có trở kháng vào nhỏ.

\given{fig5_03}{0.6\linewidth}{Opamp follower circuit\label{fig:5-03}}

Yêu cầu: chạy mô phỏng \textbf{Bias Point}, kiểm chứng $V_{OUT} = V_{(+)}$ và \textbf{thay đổi giá trị điện trở hồi tiếp $R_{15}$}.

\subsection{Mô phỏng}
\simresult
\simfigs{ex01_bias}{Kết quả Bias Point mạch follower (giá trị $R_{15}$ ban đầu)}

\begin{table}[H]
\centering
\begin{tabular}{c c c c}
\toprule
$R_{15}$ & $V_{(+)}$ (V) & $V_{(-)}$ (V) & $V_{OUT}$ (V)\\
\midrule
$1\,\text{k}\Omega$   & \fillcell & \fillcell & \fillcell\\
$10\,\text{k}\Omega$  & \fillcell & \fillcell & \fillcell\\
$100\,\text{k}\Omega$ & \fillcell & \fillcell & \fillcell\\
$1\,\text{M}\Omega$   & \fillcell & \fillcell & \fillcell\\
\bottomrule
\end{tabular}
\caption{$V_{OUT}$ theo $R_{15}$ (follower)}
\end{table}

\subsection{Tính toán}
\calc[6]
\begin{dapsobox}
\dotlineeq{$V_{OUT} =$}{với mọi $R_{15}$}
\end{dapsobox}
\remark[3]
